\chapter{Annexe A — Surface d'API}
\label{ann:api}

Toutes les routes sont servies par l'API FastAPI. Les routes de \emph{collecte}
sont soumises à la limitation de débit ; les lectures légères ne le sont pas.

\begin{table}[h]
\centering
\small
\begin{tabularx}{\textwidth}{@{}p{6.3cm}p{1.6cm}X@{}}
\toprule
\textbf{Route} & \textbf{Méthode} & \textbf{Rôle} \\
\midrule
\code{/diagnostic/adresse} & POST & Adresse \(\rightarrow\) enregistrement canonique (pipeline unitaire) \\
\code{/diagnostic/batch} & POST & Soumission d'un lot (\(N \times\) unitaire) \\
\code{/diagnostic/batch/\{id\}} & GET & Suivi d'un lot (léger, non limité) \\
\code{/diagnostic/zone/building} & GET & Fiche BDNB par identifiant de bâtiment \\
\code{/diagnostic/zone/buildings} & GET & Bâtiments d'une emprise \\
\code{/diagnostic/adresse/rapport-pdf} & GET & Rapport exportable en PDF \\
\code{/api/flood-alea} & GET & Classes de crue TRI au point (point dans polygone) \\
\code{/api/hydro} & GET & Tracé hydrographique au point (Sandre / BD TOPO) \\
\code{/api/hydro/extend} & GET & Extension amont et aval du tracé \\
\code{/api/meteo} & GET & Pluie horaire et débit prévus (Open-Meteo) \\
\code{/api/meteo/france} & GET & Synthèse nationale (moyenne de points représentatifs) \\
\code{/api/photo} & GET & Photo de terrain et crédit (Panoramax) \\
\code{/api/geocode/search} & GET & Autocomplétion d'adresse (BAN / Géoplateforme) \\
\code{/api/report} & POST & Rapport complet (gabarit + prose validée) \\
\code{/api/report/stream} & POST & Même rapport en flux SSE (patch de prose) \\
\code{/api/report/assistant} & POST & Question sur \emph{ce} rapport, avec refus explicites \\
\code{/health} & GET & Sonde de vie \\
\code{/health/detailed} & GET & Sonde détaillée \\
\bottomrule
\end{tabularx}
\caption{Surface d'API. Les routes de collecte sont énumérées \emph{exactement}
dans le limiteur de débit, ce qui évite de freiner les lectures légères.}
\end{table}

\section{Codes et états}

\begin{itemize}
  \item \textbf{200} — la réponse est produite, éventuellement
    \emph{partiellement} renseignée (dégradation déclarée dans le corps).
  \item \textbf{422} — l'entrée ne satisfait pas le schéma (question vide ou trop
    longue, champ inconnu, \code{extra="forbid"}).
  \item \textbf{429} — plafond de débit atteint sur une route de collecte.
  \item \textbf{5xx} — réservé aux erreurs réellement inattendues : une source
    publique en panne n'en produit pas, elle devient une absence déclarée.
\end{itemize}

\section{Le cas de l'assistant documentaire}

La route \code{/api/report/assistant} mérite une mention : elle ne répond qu'à
partir des faits du rapport fourni dans la requête elle-même. Aucune recherche
web, aucune mémoire d'un rapport précédent, aucune persistance côté serveur — la
conversation n'existe que dans l'échange HTTP. Une question hors périmètre
obtient un \emph{refus} (un fait sur le document) ; une défaillance du modèle
obtient une \emph{indisponibilité} (une panne de source). Les deux messages sont
distincts, et jamais remplacés par une réponse inventée.
